\section{Experimental Setup}\label{experimental-setup}

The primary experiment uses the AcademicCredentials event log from the
public Zenodo artifact associated with Chapela-Campa et al.~\cite{chapela2025quality}. The
dataset was selected because it is small enough for rapid repeated
simulation while containing the fields needed for resource-centric BPS:
case identifiers, activities, resources, start timestamps, and end
timestamps. The train and test splits contain approximately 398 cases
each, around 1,800 to 1,900 events each, 16 activities, and about 300
resources.

As a robustness dataset, the study additionally uses the LoanApp log
from the public AgentSimulator repository~\cite{kirchdorfer2024agentsimulator}. This
dataset was selected because AgentSimulator is the closest related
resource-centric BPS artifact and because the log already follows the
canonical schema required by the prototype. The prepared split contains
700 training cases and 300 test cases, with 7,492 events, 12
activities, and 19 resources. Compared with AcademicCredentials,
LoanApp has fewer resources and clearer role structure, making it a
useful test of whether resource-local and handover-aware behavior is
more visible on an AgentSimulator-style dataset.

All conditions are trained on the training split and evaluated against
the held-out test split. Each stochastic condition is repeated ten times
with different random seeds. The simulation outputs are saved as event
logs in the same canonical schema as the input logs. For the LLM-agent
proxy condition, reasoning and handover logs are also saved.

The repeated simulation protocol uses ten runs per condition. For run
\(r\) and mode index \(m\), the implementation sets the random seed to
\(1000 + 100r + m\). Each run produces a generated event log for every
condition. When full logs are saved, the LLM-agent proxy also produces
one reasoning file and one handover file per run. The repeated outputs
are aggregated into run-level and summary-level CSV/JSON files.

Two evaluation layers are used. The first layer contains lightweight
prototype metrics: trace variant distribution distance, activity
distribution distance, resource distribution distance, and mean
cycle-time relative error. These metrics are useful during development
because they are fast and interpretable. The second layer applies formal
distance measures from the Chapela-Campa et al.~framework, including
bigram and trigram control-flow distance, absolute event-time distance,
case-arrival distance, circadian event distance, workforce distance,
relative event-time distance, and cycle-time distribution distance.
Lower values indicate closer agreement with the held-out test log. The
formal repeated evaluation wraps the public
\texttt{ComputeLogDistance.py} script from the Chapela-Campa artifact
and aggregates results across the ten generated logs per condition.

